\section{Document Anatomy}
\label{sec:anatomy}

A Cooper report has exactly two page types: a first page that carries the identity block, and
running pages that carry only furniture. Understanding the split is enough to typeset a report
correctly.

The class itself is a fork of the typesetting layer used for the DeepSeek-V3 technical report
\citep{deepseek2024}, with the identity block replaced and the running header activated. The
debt is deliberate: that layout was already proven on the page sizes and print runs Cooper
cares about.

\subsection{The first page}

The header holds the wordmark on the left and the ISO-8601 date on the right
(\monthyeardate\today). Below the header rule the title is set in bold at 16\,pt, followed by the author
block and the abstract. The footer carries the copyright line at the right and, when the
corresponding author is declared, its address at the left.

The abstract is declared before \verb|\begin{document}|. This is unusual but deliberate: the
class rewrites the \texttt{abstract} environment so that its content can be placed on the title
block without floating, and the declaration has to happen before the document body starts.

\subsection{Running pages}

From page~two onwards the header is a single centred line, the short title, in 11\,pt. The
footer carries the page number, also centred. Both are separated from the text block by a 1\,pt
rule drawn in grey, which keeps the rule visible on screen without competing with body text on
paper.

\subsection{Metadata hooks}

All organisation-specific strings are collected in five macros, so that a rebrand never touches
the body of a report. Table~\ref{tab:hooks} lists them together with the class options that
switch the optional parts on.

\begin{table}[htbp]
  \centering
  \caption{The branding hooks and the class options that activate optional elements.}
  \label{tab:hooks}
  \begin{tabular}{@{}lll@{}}
    \toprule
    Hook or option & Rendered where & Default \\
    \midrule
    \verb|\reportorg|        & copyright line            & Cooper Intelligence \\
    \verb|\reportshorttitle| & running header            & Technical Report \\
    \verb|\reportlogo|       & first-page header         & none \\
    \verb|\reportaddress|    & first-page footer, left   & none \\
    \verb|\reportinternalnote| & footer when \texttt{internal} & Confidential. Internal Only. \\
    \midrule
    \texttt{logo}     & draws the wordmark               & off \\
    \texttt{copyright}& draws the copyright line         & off \\
    \texttt{address}  & draws the address line           & off \\
    \texttt{internal} & marks the report as internal     & off \\
    \texttt{nonumbering} & suppresses section numbers     & off \\
    \bottomrule
  \end{tabular}
\end{table}

Options are passed on the \verb|\documentclass| line. A released report enables \verb|logo|
and \verb|copyright|; an internal draft adds \verb|internal|, which prints the confidentiality
note in the footer.
